% !TEX program = xelatex
\documentclass[a4paper, 10pt, fontset=windows]{ctexart}

\usepackage[margin=1.8cm]{geometry}
\usepackage{booktabs}
\usepackage{enumitem}
\usepackage{xcolor}
\usepackage{array}
\usepackage{amssymb}

\setlength{\parskip}{4pt}
\setlength{\parindent}{0pt}
\pagestyle{empty}
\setlist{itemsep=1pt, leftmargin=1.5em}

\begin{document}

{\centering\Large\bfseries LabGuardian 数据采集速查\par}
\vspace{2pt}
{\centering\small\color{gray} YOLOv8-OBB $\cdot$ 面包板电路元件检测 $\cdot$ 目标 150--200 张\par}
\vspace{6pt}
\hrule
\vspace{6pt}

%% === 1 ===
\section*{1 \quad 采集目标}

\textbf{5 类元件：}
\begin{tabular}{@{}cllc@{}}
\toprule
ID & 类别 & 说明 & 引脚 \\
\midrule
0 & Resistor   & 色环电阻，水平插入，四/五环              & 2 \\
1 & Transistor & TO-92 三极管，半圆柱体                   & 3 \\
2 & Capacitor  & 电解电容(立式圆柱) / 陶瓷电容(小片)      & 2 \\
3 & Wire       & 跳线/杜邦线，多色                        & 2 \\
4 & Diode      & DO-41 二极管，黑色+银环                   & 2 \\
\bottomrule
\end{tabular}

\medskip
\textbf{6 个实验电路：}
LED限流、RC滤波、共射极放大、二极管整流、电容滤波、三极管开关(可选)

\medskip
\textbf{数据量：}
完整电路 100--120 张 + 单元件特写 20--30 张 + 空面包板 5--10 张 + 干扰场景 10--15 张 = \textbf{150--200 张}

%% === 2 ===
\section*{2 \quad 拍摄要求}

\begin{itemize}
  \item \textbf{设备}：手机后置($\geq$1200万像素，关HDR/美颜) + USB摄像头(模拟比赛)，两种都拍
  \item \textbf{分辨率}：拍最高，后处理缩放到 1280$\times$960 或 1920$\times$1080；格式 JPG
  \item \textbf{距离}：镜头距面包板 25--40cm，整个面包板在画面内
  \item \textbf{角度}：正俯拍为主(70\%)，$\pm$15$^\circ$倾斜辅助(30\%)
  \item \textbf{对焦}：点击面包板中心对焦，元件清晰即可
  \item \textbf{命名}：\texttt{exp\{编号\}\_\{光照\}\_\{序号\}.jpg}，如 \texttt{exp3\_daylight\_007.jpg}
\end{itemize}

\textbf{3 种光照（每个电路都要拍）：}
\begin{tabular}{@{}lll@{}}
L1 自然光(窗边) & L2 日光灯(教室) & L3 台灯侧光/手电筒(制造阴影)
\end{tabular}

%% === 3 ===
\section*{3 \quad 采集流程}

\textbf{准备元件：}
电阻 2--3 种(1k/10k/100$\Omega$)、三极管 2--3 个(S8050/2N3904)、电解+陶瓷电容各 2--3 个、多色杜邦线(红黑黄蓝绿各2根)、二极管 2--3 个(1N4007)。有不同面包板都拿出来。

\medskip
\textbf{拍摄矩阵（每个电路）：}

\begin{tabular}{@{}l|ccc|c@{}}
\toprule
 & L1 自然光 & L2 日光灯 & L3 侧光 & 小计 \\
\midrule
正俯拍 0$^\circ$ & 2 张 & 2 张 & 2 张 & 6 \\
左偏 15$^\circ$  & 1    & 1    & --   & 2 \\
右偏 15$^\circ$  & 1    & 1    & --   & 2 \\
\midrule
\multicolumn{4}{@{}r|}{每个电路合计} & \textbf{$\sim$10--14} \\
\bottomrule
\end{tabular}

\smallskip
6 个实验 $\times$ 14 张 $\approx$ 84 张核心图像。

\medskip
\textbf{辅助图像：}
\begin{itemize}
  \item 单元件特写(只插1--2个元件) $\sim$20张
  \item 空面包板(负样本) $\sim$5张
  \item 干扰(手在插元件、旁边有工具、元件遮挡) $\sim$10张
  \item "错误"电路(少接线、元件插反) $\sim$5张
\end{itemize}

\textbf{质检：}删掉模糊、过曝/欠曝、面包板不完整、完全重复的图片。

\vfill
{\small\color{gray}\hrule\vspace{3pt}\centering LabGuardian $\cdot$ 2026.02\par}

\end{document}
